\chapter{مفاهیم تکمیلی و سناریوهای خاص}

در گام نهایی از این مسیر آموزشی، به بررسی مباحث پراکنده و قابلیت‌های جانبی پوسته \en{Bash} می‌پردازیم. اگرچه در فصول گذشته مفاهیم بنیادین را به تفصیل شرح دادیم، اما \en{Bash} از ویژگی‌های متعددی برخوردار است که به دلیل ماهیت تخصصی‌شان، کمتر مورد بحث قرار گرفته‌اند. بسیاری از این قابلیت‌ها برای توسعه‌دهندگانی که قصد بومی‌سازی \en{Bash} در توزیع‌های لینوکس را دارند حائز اهمیت است؛ با این حال، شناخت برخی از این ابزارها برای حل چالش‌های برنامه‌نویسی خاص و پیچیده بسیار کارگشا خواهد بود. در این بخش، به تحلیل این موارد می‌پردازیم.

\section{گروه‌بندی دستورات و زیرپوسته‌ها (Subshells)}

پوسته \en{Bash} امکان گروه‌بندی دستورات را فراهم می‌آورد. این عملیات به دو روش اصلی پیاده‌سازی می‌شود: دستورات گروهی (\en{Group Commands}) و زیرپوسته‌ها (\en{Subshells}). پیش‌تر در فصل ششم، با دستورات گروهی آشنا شدیم که از نحو (\en{Syntax}) زیر پیروی می‌کنند:

\begin{latin}
\begin{lstlisting}
{ command1; command2; [command3; ...] }
\end{lstlisting}
\end{latin}

زیرپوسته‌ها نیز ساختار مشابهی دارند، با این تفاوت که از پرانتز استفاده می‌کنند:

\begin{latin}
\begin{lstlisting}
(command1; command2; [command3;...])
\end{lstlisting}
\end{latin}

تفاوت ظاهری این دو در این است که دستور گروهی با «آکولاد» و زیرپوسته با «پرانتز» محصور می‌شود. یک نکته حیاتی در مورد دستورات گروهی این است که به دلیل ساختار نحوی \en{Bash}، آکولادها حتماً باید با یک فاصله (\en{Space}) از دستورات جدا شوند و آخرین دستور نیز باید به یک سمی‌کالن (\cmd{;}) یا نویسه خط جدید (\en{Newline}) ختم گردد.

اما کاربرد عملی این دو ساختار چیست؟ فارغ از تفاوت‌های ساختاری عمیق‌تر که به آن‌ها خواهیم پرداخت، هر دو روش برای مدیریت متمرکز بازهدایت (\en{Redirection}) بسیار کارآمد هستند. فرض کنید در یک اسکریپت، قصد دارید خروجی چندین دستور را به یک فایل واحد هدایت کنید:

\begin{latin}
\begin{lstlisting}
ls -l > output.txt
echo "Listing of foo.txt" >> output.txt
cat foo.txt >> output.txt
\end{lstlisting}
\end{latin}

اگرچه روش فوق صحیح است، اما با استفاده از یک دستور گروهی می‌توان همان فرآیند را به شکلی منسجم‌تر اجرا کرد:

\begin{latin}
\begin{lstlisting}
{ ls -l; echo "Listing of foo.txt"; cat foo.txt; } > output.txt
\end{lstlisting}
\end{latin}

به همین ترتیب، استفاده از زیرپوسته نیز نتیجه مشابهی خواهد داشت:

\begin{latin}
\begin{lstlisting}
(ls -l; echo "Listing of foo.txt"; cat foo.txt) > output.txt
\end{lstlisting}
\end{latin}

مزیت اصلی این رویکرد، علاوه بر اختصار، در هنگام کار با پایپ‌لاین‌ها (\en{Pipelines}) نمایان می‌شود. ترکیب نتایج چندین دستور در یک جریان داده واحد (\en{Stream}) برای ارسال به ورودی یک دستور دیگر، با استفاده از این ساختارها بسیار ساده می‌گردد:

\begin{latin}
\begin{lstlisting}
{ ls -l; echo "Listing of foo.txt"; cat foo.txt; } | lpr
\end{lstlisting}
\end{latin}

در این مثال، خروجی هر سه دستور ادغام شده و به‌طور مستقیم به ورودی دستور \cmd{lpr} ارسال می‌شود تا یک گزارش چاپی یکپارچه تولید گردد.

علیرغم شباهت ظاهری دستورات گروهی و زیرپوسته‌ها در مدیریت جریان‌های خروجی، تفاوت عملکردی بنیادینی میان آن‌ها وجود دارد. دستورات گروهی تمامی فرامین را در درون پوسته جاری (\en{Current Shell}) اجرا می‌کنند، در حالی که زیرپوسته (همان‌طور که از نامش پیداست) یک فرآیند فرزند (\en{Child Process}) مجزا ایجاد کرده و دستورات را در آن محیط اجرا می‌کند.

هنگام ایجاد یک زیرپوسته، محیط عملیاتی آن (شامل توابع، نام‌های مستعار و متغیرها) از پوسته‌ی والد کپی و به نمونه‌ی جدید منتقل می‌شود. نکته‌ی ظریف اینجاست که زیرپوسته، برخلاف فرآیندهای فرزند معمولی، کل کپی محیط والد را به ارث می‌برد، نه فقط متغیرهای صادرشده (\en{Exported}) را. با این حال، به‌محض خاتمه‌ی زیرپوسته، این کپی از میان می‌رود و تمامی تغییراتی که در محیط آن اعمال شده‌اند (مانند مقداردهی متغیرها) نیز از بین می‌روند.

کلیدی‌ترین نکته این است که زیرپوسته‌ها برخلاف دستورات گروهی، توانایی تغییر در محیط پوسته والد را ندارند. برای درک بهتر، ابتدا عملکرد دستور گروهی را بررسی می‌کنیم:

\begin{latin}
\begin{lstlisting}
[me@linuxbox ~]$ foo="Original Value"
[me@linuxbox ~]$ { foo="Altered Value"; echo $foo; }
Altered Value
[me@linuxbox ~]$ echo $foo
Altered Value
\end{lstlisting}
\end{latin}

اکنون همین سناریو را در یک زیرپوسته تکرار می‌کنیم:

\begin{latin}
\begin{lstlisting}
[me@linuxbox ~]$ foo="Original Value"
[me@linuxbox ~]$ ( foo="Altered Value"; echo $foo )
Altered Value
[me@linuxbox ~]$ echo $foo
Original Value
\end{lstlisting}
\end{latin}

مشاهده می‌شود که دستور گروهی موفق به تغییر مقدار متغیر \cmd{foo} در محیط جاری شد، در حالی که تغییرات در زیرپوسته به همان محیط محدود ماند. این ویژگی زیرپوسته‌ها در سناریوهایی نظیر تغییر دایرکتوری موقت بسیار کارآمد است؛ چرا که تغییر مسیر با دستور \cmd{cd} در زیرپوسته، تاثیری بر دایرکتوری کاریِ پوسته والد نخواهد داشت:

\begin{latin}
\begin{lstlisting}
[me@linuxbox ~]$ pwd
/home/me
[me@linuxbox ~]$ ( cd /usr/bin; pwd )
/usr/bin
[me@linuxbox ~]$ pwd
/home/me
\end{lstlisting}
\end{latin}

به‌طور کلی، مگر در مواردی که ایزوله‌سازی فرآیند ضرورت داشته باشد، استفاده از دستورات گروهی به دلیل سرعت بالاتر و مصرف کمتر حافظه (به جهت عدم نیاز به ایجاد فرآیند جدید) ارجحیت دارد.

در اسکریپت زیر که با نام \cmd{array-2} ارائه شده، از دستورات گروهی بهره گرفته‌ایم تا تکنیک‌های پیشرفته‌تری را در کار با آرایه‌های انجمنی (\en{Associative Arrays}) نمایش دهیم. این برنامه با دریافت نام یک دایرکتوری، فهرستی از فایل‌ها را به همراه نام مالک (\en{Owner}) و گروه (\en{Group}) چاپ کرده و در انتها، آمار تجمعی فایل‌های متعلق به هر کاربر و گروه را ارائه می‌دهد.

در ادامه، نتایج خروجی اسکریپت \cmd{array-2} را (که برای اختصار تلخیص شده است) هنگام واکاوی دایرکتوری \file{/usr/bin/} مشاهده می‌کنید:

\begin{latin}
\begin{lstlisting}
[me@linuxbox ~]$ array-2 /usr/bin
/usr/bin/2to3-2.6             root    root
/usr/bin/2to3                 root    root
/usr/bin/a2p                  root    root
/usr/bin/abrowser             root    root
/usr/bin/aconnect             root    root
/usr/bin/acpi_fakekey         root    root
/usr/bin/acpi_listen          root    root
/usr/bin/add-apt-repository   root    root
.
.
/usr/bin/zipgrep              root    root
/usr/bin/zipinfo              root    root
/usr/bin/zipnote              root    root
/usr/bin/zip                  root    root
/usr/bin/zipsplit             root    root
/usr/bin/zjsdecode            root    root
/usr/bin/zsoelim              root    root

File owners:
daemon  :     1 file(s)
root    :  1394 file(s)

File group owners:
crontab :     1 file(s)
daemon  :     1 file(s)
lpadmin :     1 file(s)
mail    :     4 file(s)
mlocate :     1 file(s)
root    :  1380 file(s)
shadow  :     2 file(s)
ssh     :     1 file(s)
tty     :     2 file(s)
utmp    :   2 file(s)
\end{lstlisting}
\end{latin}

در بخش زیر، متن کامل اسکریپت به همراه شماره‌گذاری خطوط جهت تحلیل فنی ارائه شده است:

\begin{latin}
\begin{lstlisting}
 1  #!/bin/bash
 2
 3  # array-2: Use arrays to tally file owners
 4
 5  declare -A files file_group file_owner groups owners
 6
 7  if [[ ! -d "$1" ]]; then
 8      echo "Usage: array-2 dir" >&2
 9      exit 1
10  fi
11
12  for i in "$1"/*; do
13      owner="$(stat -c %U "$i")"
14      group="$(stat -c %G "$i")"
15      files["$i"]="$i"
16      file_owner["$i"]="$owner"
17      file_group["$i"]="$group"
18      ((++owners[$owner]))
19      ((++groups[$group]))
20  done
21
22  # List the collected files
23  { for i in "${files[@]}"; do
24      printf "%-40s %-10s %-10s\n" \
25          "$i" "${file_owner["$i"]}" "${file_group["$i"]}"
26  done } | sort
27  echo
28
29  # List owners
30  echo "File owners:"
31  { for i in "${!owners[@]}"; do
32      printf "%-10s: %5d file(s)\n" "$i" "${owners["$i"]}"
33  done } | sort
34  echo
35
36  # List groups
37  echo "File group owners:"
38  { for i in "${!groups[@]}"; do
39      printf "%-10s: %5d file(s)\n" "$i" "${groups["$i"]}"
40  done } | sort
\end{lstlisting}
\end{latin}

تحلیل فنی و کالبدشکافی مکانیسم اجرایی اسکریپت \cmd{array-2} به شرح زیر است:

\textbf{خط ۵:} در این بخش، آرایه‌های انجمنی با استفاده از دستور \cmd{declare} و گزینه \cmd{-A} به صورت صریح تعریف شده‌اند. در این اسکریپت، پنج آرایه با وظایف متمایز ایجاد شده‌اند:

\begin{itemize}
  \item \cmd{files}: ذخیره‌سازی نام فایل‌های دایرکتوری (با استفاده از نام فایل به عنوان کلید).
  \item \cmd{file\_group}: نگاشت مالک گروه (\en{Group Owner}) به هر فایل.
  \item \cmd{file\_owner}: نگاشت مالک (\en{User Owner}) به هر فایل.
  \item \cmd{groups}: شمارشگر تعداد فایل‌های متعلق به هر گروه.
  \item \cmd{owners}: شمارشگر تعداد فایل‌های متعلق به هر مالک.
\end{itemize}

\textbf{خطوط ۷ تا ۱۰:} این بخش شرطی اعتبار پارامتر ورودی را بررسی می‌کند. اگر مسیر ارسالی یک دایرکتوری معتبر نباشد، راهنمای استفاده از برنامه در خروجی خطای استاندارد (\en{stderr}) چاپ شده و اسکریپت با وضعیت خروجی \cmd{1} (نشان‌دهنده خطا) متوقف می‌شود.

\textbf{خطوط ۱۲ تا ۲۰:} این حلقه فرآیند پیمایش فایل‌ها را بر عهده دارد. با بهره‌گیری از دستور \cmd{stat} در خطوط ۱۳ و ۱۴، نام مالک و گروه هر فایل استخراج شده و در آرایه‌های مربوطه ذخیره می‌گردند. در خط ۱۵ نیز خودِ نام فایل به عنوان مقدار در آرایه \cmd{files} درج می‌شود تا ساختار داده‌ها تکمیل گردد.

\textbf{خطوط ۱۸ و ۱۹:} در این دو خط، از ویژگی محاسباتی پوسته برای افزایش واحد (\en{Increment}) استفاده شده است. هر بار که مالکی مشاهده می‌شود، مقدار متناظر با آن کلید در آرایه‌های \cmd{owners} و \cmd{groups} یک واحد افزایش می‌یابد تا آمار نهایی استخراج شود.

\textbf{خطوط ۲۲ تا ۲۷:} در این مرحله، فهرست فایل‌ها چاپ می‌شود. استفاده از عبارت \cmd{"\$\{files[@]\}"} تضمین می‌کند که تمامی عناصر آرایه به‌صورت کلمات مجزا پردازش شوند؛ این امر برای مدیریت صحیح فایل‌هایی که نام آن‌ها حاوی «فاصله» (\en{Space}) است، حیاتی است. نکته کلیدی در اینجا محصور کردن کل حلقه در آکولاد و تشکیل یک دستور گروهی است. از آنجا که عناصر آرایه‌های انجمنی فاقد ترتیب (\en{Order}) هستند، خروجی کل حلقه به‌صورت یکپارچه به دستور \cmd{sort} تزریق می‌شود تا فهرستی مرتب به کاربر ارائه گردد.

\textbf{خطوط ۲۹ تا ۴۰:} این دو حلقه ساختاری مشابه بخش قبل دارند، با این تفاوت که از بسط پارامتر به فرم \cmd{"\$\{!array[@]\}"} استفاده می‌کنند. این نحو نگارش به جای بازگرداندن مقادیر، فهرستی از کلیدها (اندیس‌ها) را بازمی‌گرداند که در اینجا همان نام مالکین و گروه‌ها هستند.

\section{جایگزینی فرآیند (Process Substitution)}

در فصل ۲۸، با چالش‌های محیط زیرپوسته آشنا شدیم؛ جایی که دریافتیم دستور \cmd{read} در یک پایپ‌لاین (\en{Pipeline}) مطابق انتظار عمل نمی‌کند. برای یادآوری، اگر یک پایپ‌لاین مشابه نمونه زیر ایجاد کنیم:

\begin{latin}
\begin{lstlisting}
echo "foo" | read
echo $REPLY
\end{lstlisting}
\end{latin}

مقدار متغیر \cmd{REPLY} همواره تهی خواهد بود؛ زیرا دستور \cmd{read} در یک زیرپوسته اجرا شده و متغیر \cmd{REPLY} ایجاد شده در آن، بلافاصله پس از خاتمه فرآیند زیرپوسته از بین می‌رود. از آنجا که در ساختار پایپ‌لاین، تمام دستورات در زیرپوسته‌ها اجرا می‌شوند، هر دستوری که وظیفه مقداردهی به یک متغیر را بر عهده دارد با این محدودیت مواجه خواهد شد. خوشبختانه، پوسته \en{Bash} مکانیزم خاصی به نام جایگزینی فرآیند را برای حل این معضل ارائه می‌دهد. جایگزینی فرآیند به دو صورت نحوی بیان می‌شود:

برای فرآیندهایی که خروجی استاندارد تولید می‌کنند:

\begin{latin}
\begin{lstlisting}
<(list)
\end{lstlisting}
\end{latin}

برای فرآیندهایی که ورودی استاندارد می‌پذیرند:

\begin{latin}
\begin{lstlisting}
>(list)
\end{lstlisting}
\end{latin}

در اینجا \file{list} مجموعه‌ای از دستورات اجرایی است. برای حل مشکل دستور \cmd{read} در مثال قبل، می‌توانیم از جایگزینی فرآیند به شکل زیر بهره ببریم:

\begin{latin}
\begin{lstlisting}
read < <(echo "foo")
echo $REPLY
\end{lstlisting}
\end{latin}

جایگزینی فرآیند به ما این امکان را می‌دهد که خروجی یک زیرپوسته را برای مقاصد بازهدایت (\en{Redirection})، مشابه یک «فایل معمولی» در نظر بگیریم. از آنجا که این قابلیت نوعی بسط (\en{Expansion}) محسوب می‌شود، می‌توانیم مقدار واقعی آن را مشاهده کنیم:

\begin{latin}
\begin{lstlisting}
[me@linuxbox ~]$ echo <(echo "foo")
/dev/fd/63
\end{lstlisting}
\end{latin}

با استفاده از دستور \cmd{echo} برای بررسی نتیجه بسط، درمی‌یابیم که خروجی زیرپوسته از طریق یک توصیف‌گر فایل با آدرس \file{/dev/fd/63} (یا مشابه آن) در دسترس قرار گرفته است.

این تکنیک اغلب در حلقه‌هایی که با دستور \cmd{read} کار می‌کنند استفاده می‌شود. در مثال زیر، یک حلقه \cmd{read} محتویات فهرستی از فایل‌ها را که توسط یک زیرپوسته تولید شده است، پردازش می‌کند:

\begin{latin}
\begin{lstlisting}
#!/bin/bash
# pro-sub: demo of process substitution

while read -r attr links owner group size date time filename; do
    cat << _EOF_
Filename:   $filename
Size:       $size
Owner:      $owner
Group:      $group
Modified:   $date $time
Links:      $links
Attributes: $attr

_EOF_
done < <(ls -l --time-style="+%F %H:%m"| tail -n +2)
\end{lstlisting}
\end{latin}

در این اسکریپت، حلقه برای هر سطر از فهرست فایل‌ها، دستور \cmd{read} را اجرا می‌کند. منبع داده در سطر پایانی اسکریپت تعریف شده است؛ جایی که خروجی حاصل از جایگزینی فرآیند به ورودی استاندارد (\en{stdin}) حلقه هدایت می‌شود. استفاده از دستور \cmd{tail} در پایپ‌لاین داخلی، برای حذف سطر اول خروجی \cmd{ls} (سطر \en{Total}) است که در پردازش ما مورد نیاز نیست.

پس از اجرا، اسکریپت خروجی منظمی مشابه زیر تولید خواهد کرد:

\begin{latin}
\begin{lstlisting}
[me@linuxbox ~]$ pro-sub | head -n 20
Filename:   addresses.ldif
Size:       14540
Owner:      me
Group:      me
Modified:   2009-04-02 11:12
Links:      1
Attributes: -rw-r--r--

Filename:   bin
Size:       4096
Owner:      me
Group:      me
Modified:   2009-07-10 07:31
Links:      2
Attributes: drwxr-xr-x

Filename:   bookmarks.html
Size:       394213
Owner:      me
Group:      me
\end{lstlisting}
\end{latin}

\section{پردازش پویا با دستور eval}

دستور داخلی \cmd{eval} یکی از ابزارهای چالش‌برانگیز و در عین حال قدرتمند در پوسته \en{Bash} است. در تعریف ساده، این دستور فهرستی از آرگومان‌ها را دریافت کرده، آن‌ها را به یک رشته واحد تبدیل می‌کند و سپس آن رشته را برای اجرا به خودِ پوسته بازمی‌گرداند. پرسش کلیدی این است: در چه سناریوهایی به این دستور نیاز داریم؟

در برنامه‌نویسی پوسته، گاهی لازم است یک دستور را به‌صورت پویا (\en{Dynamic})، یعنی در حین اجرای اسکریپت، تولید کنیم. بیایید با یک آزمایش ساده شروع کنیم:

می‌توان یک دستور را در قالب یک متغیر رشته‌ای ذخیره کرد و سپس با استفاده از بسط پارامتر، آن را فراخوانی نمود:

\begin{latin}
\begin{lstlisting}
[me@linuxbox ~]$ cmd="echo foo"
[me@linuxbox ~]$ $cmd
foo
[me@linuxbox ~]$
\end{lstlisting}
\end{latin}

حال بررسی می‌کنیم که اگر در دستور ما متغیر دیگری وجود داشته باشد، چه رخ می‌دهد:

\begin{latin}
\begin{lstlisting}
[me@linuxbox ~]$ str=abcde
[me@linuxbox ~]$ cmd="echo $str"
[me@linuxbox ~]$ $cmd
abcde
[me@linuxbox ~]$
\end{lstlisting}
\end{latin}

تا اینجا خروجی مطابق انتظار است. اما اگر مقدار متغیر \cmd{str} تغییر کند:

\begin{latin}
\begin{lstlisting}
[me@linuxbox ~]$ str=ABCDE
[me@linuxbox ~]$ $cmd
abcde
[me@linuxbox ~]$
\end{lstlisting}
\end{latin}

مشاهده می‌شود که خروجی تغییر نکرد؛ زیرا در لحظه مقداردهی به \cmd{cmd}، متغیر \cmd{str} بسط یافته و مقدار قبلی در آن ثابت گشته است. برای اینکه بسط متغیر را به تعویق بیندازیم، از «تک‌کوتیشن» استفاده می‌کنیم:

\begin{latin}
\begin{lstlisting}
[me@linuxbox ~]$ cmd='echo $str'
[me@linuxbox ~]$ $cmd
$str
[me@linuxbox ~]$
\end{lstlisting}
\end{latin}

در اینجا، پوسته تنها یک مرحله بسط انجام می‌دهد؛ یعنی \cmd{\$cmd} به رشته \cmd{'echo \$str'} تبدیل می‌شود، اما مرحله دوم بسط (تبدیل \cmd{\$str} به مقدار واقعی‌اش) رخ نمی‌دهد.

اینجاست که دستور \cmd{eval} وارد عمل می‌شود تا سطح دوم (یا چندم) بسط را ممکن سازد:

\begin{latin}
\begin{lstlisting}
[me@linuxbox ~]$ str=abcde
[me@linuxbox ~]$ cmd='echo $str'
[me@linuxbox ~]$ eval "$cmd"
abcde
[me@linuxbox ~]$ str=ABCDE
[me@linuxbox ~]$ eval "$cmd"
ABCDE
[me@linuxbox ~]$
\end{lstlisting}
\end{latin}

عملکرد دقیق \cmd{eval} به این صورت است: ابتدا آرگومان‌های خود را در یک رشته ادغام می‌کند، سپس روی محتوای آن رشته انواع بسط‌ها (مانند بسط تیلدا \cmd{\textasciitilde}، پارامتر \cmd{\$VAR} و جایگزینی مسیر فایل) را انجام داده و در نهایت، خروجی نهایی را در پوسته جاری (بدون ایجاد زیرپوسته) اجرا می‌کند.

\subsection{مخاطرات امنیتی دستور eval}

دستور \cmd{eval} در میان توسعه‌دهندگان به بدنامی شهرت یافته است. علت این امر، پتانسیل بالای این دستور در ایجاد حفره‌های امنیتی در اسکریپت‌های پوسته است. چنانچه رشته ورودی به \cmd{eval} از منابع خارجی (مانند ورودی‌های مستقیم کاربر) تأمین شود، ریسک اجرای فرامین مخرب بدون مجوز افزایش می‌یابد؛ پدیده‌ای که تحت عنوان «حمله تزریق کد» (\en{Code Injection}) شناخته می‌شود. بنابراین، همواره پیش از ارجاع هرگونه محتوا به \cmd{eval}، اعتبارسنجی و پالایش دقیق داده‌ها الزامی است.

\subsection{اسکریپت کمکی بازی «وردل»}

در اسکریپتی که در ادامه بررسی خواهیم کرد، از قابلیت‌های \cmd{eval} برای استخراج پاسخ‌های احتمالی در بازی «وردل» بهره می‌گیریم. برای آن دسته از مخاطبانی که با این بازی پرطرفدار آشنایی ندارند، ویکی‌پدیا آن را چنین تعریف می‌کند:

\begin{quote}
«وردل یک بازی واژگان مبتنی بر وب است که توسط مهندس نرم‌افزار ولزی، جاش واردل، طراحی و توسعه یافته است. در این بازی، بازیکنان شش فرصت برای حدس زدن یک کلمه پنج‌حرفی دارند. پس از هر حدس، بازخورد سیستم در قالب کاشی‌های رنگی ارائه می‌شود؛ این رنگ‌ها مشخص می‌کنند که کدام حروف در کلمه وجود دارند و آیا در جایگاه صحیح خود قرار گرفته‌اند یا خیر.»
\end{quote}

در اسکریپت زیر، ما از داده‌های خروجی بازی برای فیلتر کردن هوشمند فهرست کلمات و یافتن محتمل‌ترین گزینه‌ها استفاده می‌کنیم.

به طور خلاصه، پس از هر بار حدس زدن، سیستم بازی مشخص می‌کند که آیا حروف انتخاب شده در کلمه هدف وجود دارند یا خیر، و همچنین آیا در جایگاه صحیح خود قرار گرفته‌اند یا نه. بدین ترتیب، با هر تلاش اطلاعات بیشتری از واژه مجهول به دست می‌آید تا در نهایت (طبق هدف بازی) در کمترین تعداد دفعات، کلمه اصلی کشف شود.

این اسکریپت کمکی، فهرستی از واژگانی را استخراج می‌کند که با داده‌های حاصل از حدس‌های قبلی مطابقت دارند. برای این منظور، برنامه یک عبارت باقاعده (\en{Regular Expression}) ساده را دریافت می‌کند که نشان‌دهنده حروف قطعی و جایگاه آن‌هاست. علاوه بر این، فهرستی از حروف موجود (بدون جایگاه مشخص) و حروف کاملاً ناموجود در کلمه را نیز به عنوان ورودی می‌پذیرد.

فرض کنید کلمه هدف «\en{about}» باشد و ما واژه «\en{bloat}» را حدس زده باشیم. در این سناریو، بازی اعلام می‌کند که حروف «\en{o}» و «\en{t}» دقیقاً در جایگاه خود هستند، حروف «\en{a}» و «\en{b}» در کلمه وجود دارند اما در جایگاهی دیگر، و حرف «\en{l}» اصلاً در کلمه وجود ندارد. در چنین شرایطی، اسکریپت را به صورت زیر اجرا می‌کنیم:

\begin{latin}
\begin{lstlisting}
[me@linuxbox ~]$ eval-wordle ..o.t +a +b -l
\end{lstlisting}
\end{latin}

در این فرمان، آرگومان اول یک عبارت باقاعده پنج‌حرفی است که جایگاه تثبیت‌شده حروف «\en{o}» و «\en{t}» را نشان می‌دهد. سایر آرگومان‌ها نیز حروف موجود یا ناموجود را مشخص می‌کنند. خروجی اسکریپت به شرح زیر خواهد بود:

\begin{latin}
\begin{lstlisting}
abort
about
2
\end{lstlisting}
\end{latin}

این خروجی تأیید می‌کند که بر اساس معیارهای تعیین شده، تنها دو واژه در لغت‌نامه با شرایط مطابقت دارند.

این اسکریپت با پالایش محتوای یک واژه‌نامه سیستمی (\file{/usr/share/dict/words/}) بر اساس معیارهای تعیین‌شده، واژگان باقی‌مانده را استخراج می‌کند. بدین منظور، یک پایپ‌لاین (\en{Pipeline}) چندمرحله‌ای و طولانی به‌صورت پویا ایجاد می‌شود که طول آن مستقیماً به تعداد پارامترهای ورودی در خط فرمان بستگی دارد. مزیت این روش، عدم نیاز به فایل‌های موقت برای ذخیره‌سازی نتایج میانی است.

\begin{latin}
\begin{lstlisting}
 1  #!/bin/bash

 2  # eval-wordle - demonstrate eval by solving wordle puzzles

 3  PROGNAME=${0##*/}
 4  DICTIONARY=/usr/share/dict/words

 5  usage() {
 6      printf "%s\n" "Usage: ${PROGNAME} [-h|--help]"
 7      printf "%s\n" "       ${PROGNAME} regex +-char..."
 8  }

 9  help_message() {
10      cat << _EOF_
11  $PROGNAME - Wordle Helper

12  $(usage)

13  Options:
14      -h, --help     Display this help message and exit.

15  Arguments:
16  The first argument must be a five character regular
17  expression at minimum of '.....' representing five
18  unknown characters in the answer. This is followed by
19  zero or more character known to be either present or
20  absent from the answer. These are expressed as either
21  +char for letters known to be present or -char for
22  letters known to not be in the answer.

23  Example:

24  $PROGNAME a.... +o -v

25  This means we know that the answer starts with 'a' in
26  the first position and also contains an 'o' but does
27  not contain a 'v'.

  _EOF_

28   return
29  }

30  add_plus() { # Add a letter
31      local char="$1"

32      echo " | grep $char"
33      return
34  }

35  add_minus() { # Subtract a letter
36      local char="$1"

37      echo " | grep -v $char"
38      return
39  }

40  # Parse command-line

41  if [[ "$1" == "-h" || "$1" == "--help" ]]; then
42      help_message
43      exit 0
44  fi

45  if (( "${#1}" == 5 )); then
46      known_chars="$1"
47  shift
48  else
49      echo "First argument must be a 5 character regex." >&2
50      exit 1
51  fi

52  cmd="LANG=C grep '^.....$' $DICTIONARY \
53  | grep -v '[[:punct:]]' \
54  | grep -v '[[:upper:]]' \
55  | grep '$known_chars'"

56  while [[ -n "$1" ]]; do
57      case "$1" in
58          -[[:alpha:]])
59              cmd="${cmd}$(add_minus "${1:1}")"
60              ;;
61          +[[:alpha:]])
62              cmd="${cmd}$(add_plus "${1:1}")"
63              ;;
64          *)
65              echo "Invalid argument '$1'" >&2
66              exit 1
67              ;;
68      esac
69      shift
70  done

71  eval "$cmd | tee >(wc -l)"
\end{lstlisting}
\end{latin}

تحلیل ویژگی‌های فنی اسکریپت:

\textbf{خطوط ۳۰ تا ۳۹ (توابع کمکی):} دو تابع کمکی در این محدوده تعریف شده‌اند. تابع \cmd{add\_plus} (خطوط ۳۰ تا ۳۴) یک فیلتر \cmd{grep} تولید می‌کند که برای الزام حضور یک حرف خاص به کار می‌رود. تابع \cmd{add\_minus} (خطوط ۳۵ تا ۳۹) نیز یک فیلتر \cmd{grep -v} ایجاد می‌کند که برای حذف کلمات حاوی حروف غیرمجاز استفاده می‌شود.

\textbf{خطوط ۴۱ تا ۵۱ (پردازش آرگومان‌های ورودی):} در این محدوده، ابتدا خطوط ۴۱ تا ۴۴ بررسی می‌کنند که آیا آرگومان اول، درخواست راهنما (\cmd{-h} یا \cmd{--help}) است یا خیر؛ در صورت مثبت بودن، پیام راهنما نمایش داده شده و اسکریپت با کد خروج \cmd{0} متوقف می‌شود. سپس در خطوط ۴۵ تا ۵۱، از درستی آرگومان اول اطمینان حاصل می‌شود: اگر طول آن دقیقاً \cmd{5} کاراکتر نباشد، اسکریپت با پیام خطا و کد خروج \cmd{1} پایان می‌یابد.

\textbf{خطوط ۵۲ تا ۵۵ (پیکربندی اولیه پایپ‌لاین):} ساخت اولیه‌ی دستور (متغیر \cmd{cmd}) در این چهار خط انجام می‌گیرد. ابتدا (خط ۵۲) واژه‌نامه با یک عبارت باقاعده که کلمات دقیقاً ۵ حرفی را انتخاب می‌کند، پالایش می‌شود. سپس (خطوط ۵۳ و ۵۴) به‌ترتیب کلمات دارای علائم نگارشی و کلمات دارای حروف بزرگ (اسامی خاص) از نتایج حذف می‌شوند. در نهایت (خط ۵۵) فیلتر مربوط به \cmd{known\_chars}، یعنی عبارت باقاعده‌ی دریافتی از آرگومان اول، به پایپ‌لاین افزوده می‌شود.

\textbf{خط ۵۶ تا ۷۰ (حلقه‌ی پردازش آرگومان‌های اضافی):} حلقه‌ی \cmd{while} در خط ۵۶ آغاز می‌شود و ساختار \cmd{case} مربوط به آن در خط ۵۷ قرار دارد. این حلقه سایر آرگومان‌های خط فرمان را که با علامت \cmd{+} یا \cmd{-} آغاز می‌شوند پیمایش کرده (الگوهای \cmd{case} در خطوط ۵۸ و ۶۱) و با فراخوانی توابع کمکی، فیلترهای مربوطه را به متغیر \cmd{cmd} اضافه می‌کند؛ در صورت نامعتبر بودن آرگومان (خطوط ۶۴ تا ۶۷)، پیام خطا چاپ شده و اسکریپت متوقف می‌شود. حلقه در خط ۷۰ با \cmd{done} بسته می‌شود.

\textbf{خط ۷۱ (اجرای نهایی):} در این خط، دستور \cmd{eval} فرمان نهایی ذخیره‌شده در \cmd{cmd} را اجرا می‌کند. خروجی این پایپ‌لاین به دستور \cmd{tee} ارسال می‌شود تا هم در خروجی استاندارد نمایش داده شود و هم به یک «جایگزینی فرآیند» (\en{Process Substitution}) هدایت گردد که در آن دستور \cmd{wc -l} برای شمارش تعداد خطوط اجرا می‌شود؛ این روش آمار نهایی را بدون نیاز به ایجاد فایل موقت محاسبه می‌کند.

\section{مدیریت سیگنال‌ها با دستور trap}

در فصل دهم، چگونگی تعامل و پاسخ‌دهی برنامه‌ها به سیگنال‌های سیستم را بررسی کردیم. این قابلیت ارزشمند را می‌توان به اسکریپت‌های پوسته نیز تعمیم داد. اگرچه اسکریپت‌های ساده به دلیل زمان اجرای کوتاه و عدم تولید فایل‌های موقت، نیازی به این ویژگی ندارند، اما در پروژه‌های بزرگ‌تر و پیچیده‌تر، پیاده‌سازی یک روال قاعده‌مند برای مدیریت سیگنال (\en{Signal Handling}) امری ضروری است.

هنگام طراحی اسکریپت‌های جامع، باید سناریوهایی نظیر خروج ناگهانی کاربر از سیستم (\en{Logout}) یا خاموش شدن سیستم در حین اجرای برنامه را پیش‌بینی کرد. در چنین شرایطی، سیگنال‌های خاصی به تمامی فرآیندهای درگیر ارسال می‌شود. برنامه‌های استاندارد با دریافت این سیگنال‌ها، تدابیری اتخاذ می‌کنند تا خروج از برنامه به‌صورت امن و منظم انجام شود.

به‌عنوان مثال، تصور کنید اسکریپتی طراحی کرده‌اید که در طول فرآیند خود، یک فایل موقت (\en{Temporary File}) ایجاد می‌کند. بر اساس اصول طراحی بهینه، اسکریپت باید پس از تکمیل عملیات، فایل مذکور را پاکسازی کند. علاوه بر این، رویکرد هوشمندانه حکم می‌کند که اگر برنامه بر اثر دریافت یک سیگنال خارجی به‌طور غیرمنتظره متوقف شد، باز هم مکانیزمی برای حذف آن فایل و جلوگیری از انباشت داده‌های زائد وجود داشته باشد.

پوسته \en{Bash} مکانیزمی را برای این هدف در نظر گرفته است که تحت عنوان تله (\en{Trap}) شناخته می‌شود. این قابلیت از طریق فرمان داخلی \cmd{trap} پیاده‌سازی شده و از نحو (\en{Syntax}) زیر پیروی می‌کند:

\begin{latin}
\begin{lstlisting}
trap argument signal [signal...]
\end{lstlisting}
\end{latin}

در این ساختار، \file{argument} رشته‌ای حاوی فرامینی است که باید اجرا شوند و \file{signal} معرف سیگنالی است که ماشه اجرای آن فرامین را فعال می‌کند.

در ادامه، یک نمونه ساده از مدیریت سیگنال ارائه شده است:

\begin{latin}
\begin{lstlisting}
#!/bin/bash
# trap-demo: simple signal handling demo

trap "echo 'I am ignoring you.'" SIGINT SIGTERM

for i in {1..5}; do
    echo "Iteration $i of 5"
    sleep 5
done
\end{lstlisting}
\end{latin}

این اسکریپت تله‌ای را تعریف می‌کند که با دریافت سیگنال‌های \en{SIGINT} (ارسال شده توسط کلیدهای ترکیبی \cmd{Ctrl+c}) یا \en{SIGTERM} (سیگنال خاتمه استاندارد)، دستور \cmd{echo} را اجرا می‌کند. خروجی برنامه هنگام تلاش کاربر برای متوقف کردن آن به شرح زیر است:

\begin{latin}
\begin{lstlisting}
[me@linuxbox ~]$ trap-demo
Iteration 1 of 5
Iteration 2 of 5
^CI am ignoring you.
Iteration 3 of 5
^CI am ignoring you.
Iteration 4 of 5
Iteration 5 of 5
\end{lstlisting}
\end{latin}

همان‌طور که مشاهده می‌شود، با هر بار تلاش کاربر برای قطع برنامه، پیام جایگزین چاپ شده و اسکریپت به کار خود ادامه می‌دهد.

از آنجا که گنجاندن توالی پیچیده‌ای از دستورات در قالب یک رشته متنی دشوار و غیراصولی است، روش معمول آن است که یک تابع (\en{Function}) برای مدیریت سیگنال تعریف و فراخوانی شود. در مثال زیر، برای هر سیگنال ورودی، یک تابع مجزا در نظر گرفته شده است:

\begin{latin}
\begin{lstlisting}
#!/bin/bash

# trap-demo2: simple signal handling demo

exit_on_signal_SIGINT () {
    echo "Script interrupted." 2>&1
    exit 0
}

exit_on_signal_SIGTERM () {
    echo "Script terminated." 2>&1
    exit 0
}

trap exit_on_signal_SIGINT SIGINT
trap exit_on_signal_SIGTERM SIGTERM

for i in {1..5}; do
    echo "Iteration $i of 5"
    sleep 5
done
\end{lstlisting}
\end{latin}

این اسکریپت از دو دستور \cmd{trap} بهره می‌برد که هر کدام به تابعی خاص متصل شده‌اند. نکته حائز اهمیت، وجود دستور \cmd{exit} در انتهای این توابع است؛ چرا که بدون فراخوانی صریح \cmd{exit}، اسکریپت پس از اجرای بدنهٔ تله، به روند عادی خود بازگشته و متوقف نخواهد شد.

هنگامی که کاربر در حین اجرای این اسکریپت از کلید ترکیبی \cmd{Ctrl+c} استفاده می‌کند، خروجی به شرح زیر خواهد بود:

\begin{latin}
\begin{lstlisting}
[me@linuxbox ~]$ trap-demo2
Iteration 1 of 5
Iteration 2 of 5
^CScript interrupted.
\end{lstlisting}
\end{latin}

فرمان داخلی \cmd{trap} علاوه بر سیگنال‌های استاندارد سیستم‌عامل، از چندین سیگنال اختصاصی پوسته (\en{Internal Shell Signals}) نیز پشتیبانی می‌کند که در مدیریت خطاهای اجرایی بسیار کارآمد هستند. دو مورد از کاربردی‌ترینِ این سیگنال‌ها، \en{EXIT} و \en{ERR} نام دارند.

تله \en{EXIT} به محض خاتمه یافتن اسکریپت (به هر دلیلی) فعال می‌شود. تله \en{ERR} نیز زمانی فعال می‌شود که یک فرمان با وضعیت خروجی (\en{Exit Status}) غیرصفر پایان یابد. این قابلیت، در واقع نسخه‌ای هوشمندتر و انعطاف‌پذیرتر از گزینه \cmd{set -e} است که پیش‌تر در مبحث عیب‌یابی اسکریپت‌ها بررسی کردیم.

در ادامه، اسکریپتی برای نمایش عملکرد این دو تله ارائه شده است. به خط شماره ۵ توجه کنید که حاوی یک خطای عمدی برای تحریک تله است:

\begin{latin}
\begin{lstlisting}
 1  #!/bin/bash

 2  # trap-demo3 - demonstrate ERR and EXIT signal handling

 3  trap "echo \"There is an error.\"" ERR
 4  trap "echo \"The program has ended.\"" EXIT

 5  echox "Running..."

 6  read -r -p "Say something... " something
 7  echo "$something"
\end{lstlisting}
\end{latin}

در این اسکریپت، ابتدا تله‌ها تعریف می‌شوند. سپس در خط ۵، سعی می‌شود پیامی چاپ شود، اما دستور \cmd{echo} به اشتباه \cmd{echox} تایپ شده است. به محض بروز این خطا، تله \en{ERR} فعال شده و پیام خطا را چاپ می‌کند. با این حال، برخلاف دستور \cmd{set -e}، اسکریپت متوقف نمی‌شود و به خط بعدی می‌رود تا منتظر ورودی کاربر بماند. در نهایت، پس از اتمام اجرای آخرین دستور، تله \en{EXIT} فعال شده و پیام پایان برنامه را نمایش می‌دهد.

پس از اجرای اسکریپت فوق، خروجی عملیاتی به شرح زیر در ترمینال ظاهر می‌شود:

\begin{latin}
\begin{lstlisting}
[me@linuxbox ~]$ trap-demo3
./trap-demo3: line 8: echox: command not found
There is an error.
Say something...

The program has ended.
\end{lstlisting}
\end{latin}

تحلیل توالی رویدادها در این خروجی بسیار آموزنده است. ابتدا پیام خطای استاندارد پوسته مشاهده می‌شود که از عدم وجود دستور \cmd{echox} خبر می‌دهد. بلافاصله پس از آن، تله \en{ERR} فعال شده و رشته تعیین‌شده را در خروجی چاپ می‌کند. در گام بعد، اسکریپت متوقف نشده و دستور \cmd{read} اجرا می‌شود. با فشردن کلید \en{ENTER} توسط کاربر، فرآیند اجرای دستورات به پایان رسیده و در نهایت تله \en{EXIT} فعال می‌گردد تا اتمام برنامه را اعلام کند.

این رفتار نشان‌دهنده تفاوت بنیادین میان مدیریت خطا با استفاده از تله‌ها و توقف اجباری اسکریپت است؛ در واقع تله‌ها به ما اجازه می‌دهند بدون قطع لزوماً روند برنامه، نسبت به رخدادهای غیرمنتظره واکنش نشان دهیم.

\section{مدیریت فایل‌های موقت (Temporary Files)}

یکی از دلایل اصلی تعبیه «مدیریت‌کننده‌های سیگنال» (\en{Signal Handlers}) در اسکریپت‌ها، پاک‌سازی فایل‌های موقتی است که برنامه برای ذخیره نتایج میانی در طول چرخه حیات خود ایجاد می‌کند. نام‌گذاری اصولی فایل‌های موقت در دنیای یونیکس و لینوکس فراتر از یک انتخاب ساده و در واقع نوعی مهارت فنی است. به‌طور سنتی، برنامه‌ها در سیستم‌های شبه‌یونیکس فایل‌های موقت خود را در مسیر \file{/tmp/} ایجاد می‌کنند؛ دایرکتوری مشترکی که دقیقاً برای همین منظور طراحی شده است. با این حال، ماهیت اشتراکی این دایرکتوری مخاطرات امنیتی خاصی را به‌ویژه برای اسکریپت‌هایی که با دسترسی ابرکاربر (\en{Root}) اجرا می‌شوند، به همراه دارد.

علاوه بر گام بدیهیِ تنظیم سطوح دسترسی (\en{Permissions}) مناسب برای فایل‌هایی که در محیطی عمومی قرار دارند، «پیش‌بینی‌ناپذیر بودن» نام این فایل‌ها اهمیتی حیاتی دارد. این رویکرد مانع از نوعی آسیب‌پذیری امنیتی به نام شرایط رقابتی در فایل‌های موقت (\en{Race Condition in Temp Files}) می‌شود. یکی از روش‌های سنتی برای ایجاد نامی پیش‌بینی‌ناپذیر اما توصیفی، استفاده از الگوی زیر است:

\begin{latin}
\begin{lstlisting}
tempfile=/tmp/$(basename $0).$$.$RANDOM
\end{lstlisting}
\end{latin}

این دستور نام فایلی را می‌سازد که ترکیبی از نام اسکریپت، شناسه فرآیند (\en{PID}) و یک عدد تصادفی است. با این حال، باید توجه داشت که متغیر \cmd{\$RANDOM} در پوسته \en{Bash} تنها عددی بین ۱ تا ۳۲۷۶۷ تولید می‌کند؛ این بازه در مقیاس محاسباتی چندان بزرگ نیست و استفاده از آن به‌تنهایی ممکن است در برابر یک مهاجم حرفه‌ای و مصمم کارساز نباشد.

روش استاندارد و ایمن‌تر، بهره‌گیری از ابزار \cmd{mktemp} است (که نباید با تابع کتابخانه‌ای \cmd{mktemp} در زبان \en{C} اشتباه گرفته شود). این برنامه یک «الگو» (\en{Template}) را به عنوان آرگومان دریافت کرده و بر اساس آن، فایل موقت را ایجاد و نام‌گذاری می‌کند. الگو باید شامل رشته‌ای از کاراکترهای \cmd{"X"} باشد که با حروف و اعداد تصادفی جایگزین می‌شوند. هرچه تعداد \cmd{"X"}ها بیشتر باشد، پیچیدگی و امنیت نام تولید شده بالاتر خواهد بود:

\begin{latin}
\begin{lstlisting}
tempfile=$(mktemp /tmp/foobar.$$.XXXXXXXXXX)
\end{lstlisting}
\end{latin}

این دستور فایل را ایجاد کرده و مسیر دقیق آن را در متغیر \cmd{tempfile} ذخیره می‌کند. کاراکترهای الگو با مقادیر تصادفی جایگزین شده و خروجی نهایی (که در این مثال \en{PID} اسکریپت را نیز شامل می‌شود) چیزی شبیه به نمونه زیر خواهد بود:

\begin{latin}
\begin{lstlisting}
/tmp/foobar.6593.UOZuvM6654
\end{lstlisting}
\end{latin}

برای اسکریپت‌هایی که توسط کاربران عادی اجرا می‌شوند، به‌جای استفاده از دایرکتوری عمومی \file{/tmp/}، توصیه می‌شود یک دایرکتوری اختصاصی برای فایل‌های موقت در دایرکتوری خانگی کاربر ایجاد شود، به‌عنوان مثال:

\begin{latin}
\begin{lstlisting}
[[ -d $HOME/tmp ]] || mkdir $HOME/tmp
\end{lstlisting}
\end{latin}

\section{اجرای ناهمگام (Asynchronous Execution)}

در برخی سناریوها، اجرای هم‌زمان چندین وظیفه به‌جای پردازش ترتیبی آن‌ها ضرورت می‌یابد. همان‌طور که می‌دانیم، تمامی سیستم‌عامل‌های مدرن از قابلیت چندوظیفگی (\en{Multitasking}) و در بسیاری از موارد چندکاربری (\en{Multiuser}) پشتیبانی می‌کنند. اسکریپت‌های پوسته نیز می‌توانند به‌گونه‌ای طراحی شوند که از این ظرفیت برای اجرای موازی عملیات بهره ببرند.

این سازوکار معمولاً شامل اجرای یک اسکریپت والد است که یک یا چند اسکریپت فرزند را برای انجام وظایف جانبی فراخوانی می‌کند، در حالی که خود به اجرای دستورات اصلی ادامه می‌دهد. با این حال، زمانی که مجموعه‌ای از اسکریپت‌ها به‌صورت ناهمگام اجرا می‌شوند، چالش هماهنگ‌سازی (\en{Synchronization}) بین فرآیند والد و فرزند پدید می‌آید. به‌عبارت دیگر، اگر اسکریپت والد برای تکمیل نهایی کار خود به نتایج اسکریپت فرزند وابسته باشد، باید مکانیزمی برای انتظار و تایید اتمام کار فرزند وجود داشته باشد.

پوسته \en{Bash} برای مدیریت این تعاملات ناهمگام، فرمان داخلی \cmd{wait} را ارائه کرده است. استفاده از دستور \cmd{wait} باعث می‌شود که اجرای اسکریپت والد تا زمان خاتمه یافتن یک فرآیند مشخص (معمولاً اسکریپت فرزند که در پس‌زمینه اجرا شده است) متوقف بماند.

\subsection{هماهنگ‌سازی فرآیندها با دستور wait}

برای درک بهتر عملکرد دستور \cmd{wait} در مدیریت فرآیندهای ناهمگام، از یک سناریوی عملی شامل دو اسکریپت «والد» و «فرزند» استفاده می‌کنیم. ابتدا اسکریپت والد را بررسی می‌کنیم که وظیفه مدیریت چرخه حیات فرآیند فرزند را بر عهده دارد:

\begin{latin}
\begin{lstlisting}
#!/bin/bash

# async-parent: Asynchronous execution demo (parent)

echo "Parent: starting..."

echo "Parent: launching child script..."
async-child &
pid=$!

echo "Parent: child (PID= $pid) launched."

echo "Parent: continuing..."
sleep 2

echo "Parent: pausing to wait for child to finish..."
wait "$pid"

echo "Parent: child is finished. Continuing..."
echo "Parent: parent is done. Exiting."
\end{lstlisting}
\end{latin}

در مرحله بعد، اسکریپت فرزند را تعریف می‌کنیم که یک وظیفه زمان‌بر (در اینجا یک توقف ۵ ثانیه‌ای) را شبیه‌سازی می‌کند:

\begin{latin}
\begin{lstlisting}
#!/bin/bash

# async-child: Asynchronous execution demo (child)

echo "Child: child is running..."
sleep 5
echo "Child: child is done. Exiting."
\end{lstlisting}
\end{latin}

در این مدل، اسکریپت فرزند ساختار ساده‌ای دارد و منطق اصلی در اسکریپت والد پیاده‌سازی شده است. فرآیند والد پس از اجرای فرآیند فرزند، آن را با استفاده از عملگر \cmd{\&} به پس‌زمینه (\en{Background}) می‌فرستد. بلافاصله با بهره‌گیری از پارامتر ویژه \cmd{\$!} (که همواره حاوی \en{PID} آخرین پردازش اجرا شده در پس‌زمینه است)، شناسه فرآیند فرزند در متغیر \cmd{pid} ذخیره می‌شود.

اسکریپت والد پس از انجام برخی کارهای موازی، به نقطه‌ای می‌رسد که برای ادامه مسیر به اتمام کار فرآیند فرزند نیاز دارد. در این مرحله، دستور \cmd{wait} به همراه \en{PID} فرزند فراخوانی می‌شود. این دستور باعث می‌شود تا اجرای اسکریپت والد دقیقاً در همین خط متوقف (\en{Block}) شود و تنها زمانی ادامه یابد که فرآیند فرزند با موفقیت یا خطا خاتمه یافته و سیگنال خروج خود را ارسال کند.

پس از اجرای هم‌زمان، خروجی تعامل میان اسکریپت‌های والد و فرزند به شرح زیر در ترمینال نمایش داده می‌شود:

\begin{latin}
\begin{lstlisting}
[me@linuxbox ~]$ async-parent
Parent: starting...
Parent: launching child script...
Parent: child (PID= 6741) launched.
Parent: continuing...
Child: child is running...
Parent: pausing to wait for child to finish...
Child: child is done. Exiting.
\end{lstlisting}
\end{latin}

\section{پایپ‌های نام‌گذاری‌شده (Named Pipes)}

در اکثر سیستم‌های شبه‌یونیکس، امکان ایجاد نوع خاصی از فایل تحت عنوان پایپ نامگذاری شده (\en{Named Pipes}) وجود دارد. این فایل‌ها ابزاری برای برقراری ارتباط میان دو فرآیند مستقل هستند و می‌توان از آن‌ها مشابه فایل‌های معمولی در سیستم فایل استفاده کرد. اگرچه امروزه این رویکرد کمتر در اولویت توسعه‌دهندگان است، اما درک مکانیسم آن برای مدیریت سیستم‌های لینوکسی بسیار حائز اهمیت است.

در مهندسی نرم‌افزار، معماری رایجی به نام کلاینت-سرور (\en{Client-Server}) وجود دارد که می‌تواند از روش‌های مختلف ارتباط بین‌فرآیندی (\en{IPC}) نظیر پایپ نامگذاری شده یا اتصالات شبکه بهره ببرد. ملموس‌ترین نمونه از این معماری، تعامل میان مرورگر وب (در نقش کلاینت) و سرور وب است؛ کلاینت درخواست‌های خود را ارسال کرده و سرور با داده‌های مورد نظر (صفحات وب) پاسخ می‌دهد.

پایپ نامگذاری شده در سطح سیستم فایل مانند یک فایل عادی ظاهر می‌شوند، اما در واقع بافرهایی از نوع \en{FIFO} (اولین ورودی، اولین خروجی) هستند. مشابه پایپ‌لاین‌های معمولی (\en{Pipes})، داده‌ها از یک سمت وارد شده و با همان ترتیب از سمت دیگر خارج می‌شوند. با استفاده از این پایپ‌ها، می‌توان ساختاری مشابه زیر ایجاد کرد:

\begin{latin}
\begin{lstlisting}
process1 > named_pipe
\end{lstlisting}
\end{latin}

و در سمتی دیگر:

\begin{latin}
\begin{lstlisting}
process2 < named_pipe
\end{lstlisting}
\end{latin}

این ساختار در عمل رفتاری مشابه پایپ‌لاین استاندارد (\cmd{process1 | process2}) دارد، با این تفاوت که در اینجا ارتباط از طریق یک گره (\en{Node}) مشخص در سیستم فایل برقرار می‌گردد.

\subsection{ایجاد پایپ نام‌گذاری‌شده (Named Pipe)}

برای شروع، ابتدا باید یک پایپ نامگذاری شده ایجاد کنیم. این عملیات با استفاده از فرمان \cmd{mkfifo} به سادگی قابل انجام است:

\begin{latin}
\begin{lstlisting}
[me@linuxbox ~]$ mkfifo pipe1
[me@linuxbox ~]$ ls -l pipe1
prw-r--r-- 1 me me 0 2009-07-17 06:41 pipe1
\end{lstlisting}
\end{latin}

در مثال فوق، پایپی به نام \file{pipe1} ایجاد شده است. با بررسی خروجی دستور \cmd{ls -l}، مشاهده می‌کنید که اولین کاراکتر در بخش مشخصه‌های فایل (\en{File Attributes})، حرف «\cmd{p}» است؛ این علامت در سیستم‌عامل‌های شبه‌یونیکس به معنای \en{Pipe} بوده و نشان می‌دهد که این فایل یک پایپ نامگذاری شده (\en{FIFO}) است و نه یک فایل معمولی یا دایرکتوری.

\subsection{تعامل با پایپ نام‌گذاری‌شده}

برای درک سازوکار عملیاتی یک پایپ نامگذاری شده، از دو پنجره ترمینال (یا دو کنسول مجازی مجزا) استفاده می‌کنیم. در ترمینال اول، یک دستور ساده را اجرا کرده و خروجی آن را به سمت مجرا هدایت می‌کنیم:

\begin{latin}
\begin{lstlisting}
[me@linuxbox ~]$ ls -l > pipe1

\end{lstlisting}
\end{latin}

پس از فشردن کلید \en{Enter}، مشاهده می‌شود که خط فرمان در حالت تعلیق باقی می‌ماند و به ظاهر «قفل» شده است. این رخداد به این دلیل است که در سمت دیگر مجرا، هنوز هیچ فرآیندی برای دریافت و خواندن داده‌ها متصل نشده است. در ادبیات سیستم‌عامل، به این وضعیت مسدود شدن (\en{Blocking}) مجرا گفته می‌شود. این توقف تا زمانی ادامه می‌یابد که فرآیند دیگری به انتهای دیگر مجرا متصل شده و عملیات خواندن داده را آغاز کند.

اکنون در پنجره دوم ترمینال، دستور زیر را وارد می‌کنیم:

\begin{latin}
\begin{lstlisting}
[me@linuxbox ~]$ cat < pipe1
\end{lstlisting}
\end{latin}

با اجرای این دستور، فهرست فایل‌هایی که توسط فرمان \cmd{ls} در ترمینال اول تولید شده بود، در خروجی دستور \cmd{cat} در ترمینال دوم ظاهر می‌شود. به محض برقراری این اتصال و تخلیه بافر، فرآیند \cmd{ls} در ترمینال اول از حالت مسدود خارج شده و با موفقیت به پایان می‌رسد.

\section{جمع‌بندی}

بدین ترتیب به پایان این مسیر آموزشی می‌رسیم. از این نقطه به بعد، تداوم یادگیری تنها در گرو «تمرین، تمرین و تکرار» است. اگرچه مباحث گسترده‌ای را با هم مرور کردیم، اما دنیای لینوکس همچنان قلمروهای کشف‌نشده بسیاری دارد.

هزاران ابزار در محیط خط فرمان (\en{CLI}) منتظر بهره‌برداری هوشمندانه شما هستند. پیشنهاد می‌شود جست‌وجو در مسیر \file{/usr/bin/} را آغاز کنید تا با تنوع خیره‌کننده برنامه‌های موجود آشنا شوید. برای آن دسته از مخاطبانی که مشتاق یادگیری عمیق‌تر هستند، کتاب تکمیلی نویسنده‌ی اصلی با عنوان «ماجراجویی با خط فرمان لینوکس» (\en{Adventures with the Linux Command Line}) به‌صورت رایگان در وب‌سایت \en{LinuxCommand.org} برای دانلود در دسترس است.

\section{منابع مطالعه تکمیلی}

بخش «دستورات مرکب» (\en{Compound Commands}) در صفحات راهنمای \en{Bash}، شرح کاملی از نشانه‌گذاری‌های دستورات گروهی (\en{group commands}) و زیرپوسته‌ها (\en{subshells}) ارائه می‌دهد.

بخش بسط‌ها (\en{EXPANSION}) در صفحات راهنمای \en{Bash}، شامل بندی است که به مبحث جایگزینی فرآیند (\en{process substitution}) می‌پردازد.

راهنمای پیشرفته اسکریپت‌نویسی \en{Bash} (\en{Advanced Bash-Scripting Guide}) نیز بحثی پیرامون جایگزینی فرآیند ارائه کرده است:

\begin{latin}
\url{http://tldp.org/LDP/abs/html/process-sub.html}
\end{latin}

مجله لینوکس (\en{Linux Journal}) دو مقاله مفید درباره پایپ‌های نام‌گذاری‌شده (\en{named pipes}) دارد. مقاله نخست، منتشرشده در سپتامبر ۱۹۹۷:

\begin{latin}
\url{http://www.linuxjournal.com/article/2156}
\end{latin}

و مقاله دوم، منتشرشده در مارس ۲۰۰۹:

\begin{latin}
\url{http://www.linuxjournal.com/content/using-named-pipes-fifos-bash}
\end{latin}